\section{Ethics and responsible use}\label{sec:ethics}

Models that forecast surges in hate incidents can protect people. They can
also cause harm through how outputs are worded, where they are shared and
what they are used for. This section sets out the protocol that accompanies
the plan in \cref{sec:plan}. It complements the companion paper's protocol
for victim privacy and record linkage \citep{kornsteinlab2026darkfigure}.

\subsection{Risks of forecasting}\label{sec:ethics-risks}

\paragraph{Stigmatisation.}
A map of elevated predicted intensity can be read as a statement about the
people who live in an area. That reading applies to neighbourhoods where
Jewish residents are targeted, which can come to be seen as ``dangerous''.
It applies equally to any community that observers might wrongly treat as a
source of incidents. The model contains no information about perpetrators,
and its spatial units are recording areas, not communities. Alerts and published maps are therefore
at the borough level or coarser, and never street- or venue-level heat maps.
Finer units used in sensitivity analysis S6 are reported only as summary
statistics. No output ranks neighbourhoods or communities.

\paragraph{Self-fulfilling and self-defeating alerts.}
An alert can change what it predicts. Publicity can raise reporting, which
inflates recorded counts (the time-varying $p_k$ of analysis S2). Protective
deployment can prevent incidents, which makes a correct alert look like a
false alarm. Either effect contaminates prospective evaluation. Alerts
during the prospective period (phase P7) are therefore shared only under the
protocol below, and evaluation reports state which alerts were acted on.

\paragraph{Information hazards.}
Forecasts that name ``hot'' dates or locations could help someone planning
an attack, and could encourage imitation. We therefore publish methods,
code and retrospective results, but no live or near-live forecasts and no
forecast tied to a specific site. Retrospective results are released only
after an embargo (\cref{sec:ethics-disclosure}).

\subsection{Responsible disclosure of alerts}\label{sec:ethics-disclosure}

Alerts produced during the prospective evaluation are handled as follows.
\begin{enumerate}[1.]
  \item \emph{Recipients.} Alerts go only to vetted community-security
  partners, and to a police liaison where the partner agreement provides for
  this. Each partner signs an agreement that limits use to protective
  purposes, such as security staffing and victim support, and excludes
  enforcement targeting of any community.
  \item \emph{Content.} An alert gives the borough or coarser area, the
  window $\Delta$, the predicted probability of exceeding the threshold
  with its uncertainty, and a fixed statement of the model's limitations.
  It contains no individual-level information and no attribution to
  persons, groups or communities.
  \item \emph{Human review.} A named member of the research team reviews
  every alert before release, and may withhold an alert that is plainly
  driven by a data artefact.
  \item \emph{Channel and logging.} Alerts are sent over an encrypted
  channel. Each alert, the decision to release or withhold it, and the
  partner's acknowledgement are logged for audit.
  \item \emph{Embargo.} Alerts are not published. Aggregate evaluation
  results (hit and false-alarm rates by quarter and borough) are published
  no earlier than 90 days after the end of the period they cover.
  \item \emph{Suspension.} Alerts are suspended if rolling calibration
  deteriorates beyond a pre-specified tolerance, or if a partner reports
  misuse.
\end{enumerate}

\subsection{Language and attribution}\label{sec:ethics-language}

The model attributes \emph{timing}, not blame. A positive shock loading
$\beta_k$, or a large direct component $D_k$ after event $e$, means that
recorded incidents in stream $k$ rose after $\tau_e$ by more than the
background and the contagion terms predict. It does not mean that the
parties to the event, the communities associated with them, or the people
who discussed it caused any incident. Incidents are committed by
individual perpetrators, about whom the model contains no information. We
therefore:
\begin{enumerate}[(i)]
  \item describe events as \emph{salient external events} of a stated class
  and date, and do not characterise them politically;
  \item never describe a community, nationality or religious group as a
  source, driver or cause of incidents;
  \item phrase every result in terms of recorded incidents and statistical
  association; this includes the online$\to$offline coefficients, which
  are Granger-type (\cref{prop:granger}); and
  \item report event-level decompositions in an appendix under neutral
  identifiers, with class-level summaries in the main text.
\end{enumerate}

\subsection{Political neutrality}\label{sec:ethics-neutrality}

The paper takes no position on any political conflict, government, movement
or public figure. Neutrality is built into the procedure, not just asserted:
event classes are defined by type and apply to all parties alike
(\cref{sec:plan-events}); the event list is frozen before outcome data are
seen, so events cannot be selected after the fact; every pre-registered
event is reported whatever the result (rule D3); and the event-coding rules
and frozen list are public. Before registration, the event-coding protocol
will be reviewed by at least two external reviewers with differing
perspectives on the events likely to be coded. Their comments and our
responses will be deposited with the registration.

\subsection{Online data}\label{sec:ethics-online}

Online data are used only under the terms of the access route through which
they are obtained, including any conditions on data protection and security
attached to access under Article~40 of the Digital Services Act
\citep{eu2022dsa}. The analysis uses aggregate counts per time bin. We do
not build account-level scores or profiles, do not link online accounts to
offline incidents, victims or suspects, and do not publish post text,
handles or quotations that would allow re-identification. Annotation
samples for estimating classifier error (\cref{eq:data-classerr}) are held
under access control and deleted when the retention period agreed with the
data provider and the ethics board ends. Classifier outputs are used only as
aggregate series. They are never used to flag individuals.

\subsection{Ethics review and data security}\label{sec:ethics-irb}

The protocol, including this section, will be submitted to an institutional
review board or equivalent research ethics committee before any data are
accessed (gate F5). This applies even to publicly downloadable data, because
the plan combines several streams and may involve partner and platform
data. Restricted data are kept on encrypted storage with role-based access,
and use is logged. The custodian who holds the incident data during event
coding (\cref{sec:plan-events}) is separate from the coders. Outputs follow
a minimum-cell rule: counts below 10 are not shown below the borough level,
and any table that could reveal a single incident's date and location is
suppressed or coarsened. Derived counts are released at a resolution no
finer than the source itself publishes; for restricted streams, they are
released weekly or more coarsely. Data are retained no longer than the agreements
allow, and the replication package contains only data whose release is
permitted (\cref{sec:plan-replication}).
